\documentclass[11pt]{article}
\usepackage[utf8]{inputenc}
\usepackage[T1]{fontenc}
\usepackage{lmodern}
\usepackage[spanish]{babel}
\usepackage[margin=2.5cm]{geometry}
\usepackage{amsmath,amssymb}
\usepackage{parskip}
\usepackage{xcolor}
\usepackage[colorlinks=true,linkcolor=black,urlcolor=blue,citecolor=black]{hyperref}

\title{\textbf{Análisis de sentimientos sobre el Codex Gigas:\\propuesta de investigación}}
\author{Antonio Soria\\ \small{Divulgación Científica — purpplealien.github.io}}
\date{Julio 2026}

\begin{document}
\maketitle

\begin{abstract}
El Codex Gigas --- la ``Biblia del Diablo'' --- es el manuscrito medieval más grande que se
conserva: 75 kilogramos de vitela escritos, según el análisis paleográfico, por una sola mano
a lo largo de décadas. Este artículo presenta mi próxima línea de investigación: aplicar
procesamiento de lenguaje natural y modelos de lenguaje contextuales para latín al contenido
del códice, con el objetivo de construir un \emph{mapa emocional} computacional de sus textos,
de los evangelios a las fórmulas de exorcismo. Se describen el corpus, la metodología
propuesta, los retos técnicos y los resultados esperados.
\end{abstract}

\section{El objeto: un libro imposible}

El Codex Gigas fue creado a inicios del siglo XIII (se estima entre 1204 y 1230) en el
monasterio benedictino de Podlažice, en Bohemia --- la actual República Checa. Sus números
intimidan: 92~cm de alto, 50 de ancho, 22 de grosor, 74.8~kg de peso y 310 hojas de vitela
que habrían requerido las pieles de más de 150 animales. El análisis de la caligrafía indica
un solo escriba --- la tradición lo llama Herman el Recluso --- trabajando quizá veinte o
treinta años.

La leyenda que le da su apodo cuenta que un monje condenado a ser emparedado vivo prometió
escribir en una sola noche un libro con todo el conocimiento humano, y al verse perdido cerca
de la medianoche vendió su alma al diablo para terminarlo. Como firma del pacto, el folio 290
recto exhibe un retrato del demonio a página completa --- exactamente frente a una
representación de la Ciudad Celestial. Desde 1648, cuando las tropas suecas lo tomaron como
botín de guerra en Praga, se conserva en la Biblioteca Nacional de Suecia en Estocolmo, que
mantiene una digitalización completa de acceso público.

\subsection{El contenido: una biblioteca en un solo lomo}

Para un proyecto de NLP, lo valioso del códice es su heterogeneidad textual. Contiene, en
latín: la Biblia completa (Antiguo y Nuevo Testamento en versión Vulgata), las
\emph{Antigüedades} y \emph{La guerra de los judíos} de Flavio Josefo, las \emph{Etimologías}
de Isidoro de Sevilla (la enciclopedia estándar medieval), el \emph{Ars medicinae} y textos
médicos de la escuela salernitana, la \emph{Chronica Boemorum} de Cosmas de Praga, un
calendario con necrologio, listas de nombres, y --- lo más jugoso para un análisis emocional
--- conjuros mágicos y fórmulas de exorcismo. Géneros radicalmente distintos: narrativa
sagrada, historia, enciclopedia, medicina, crónica local y magia ritual, todos copiados por
la misma mano.

\section{La pregunta de investigación}

El análisis de sentimientos (\emph{sentiment analysis}) detecta computacionalmente la
polaridad (positivo/negativo) y las emociones (miedo, alegría, ira\ldots) expresadas en un
texto. Aplicado al Codex Gigas, propongo responder tres preguntas:

\begin{enumerate}
  \item ¿Cómo se distribuye el contenido emocional a lo largo del códice? ¿Existe un
        \emph{perfil afectivo} distinguible por género textual --- evangelios frente a
        crónicas frente a exorcismos?
  \item ¿Las fórmulas de exorcismo y los conjuros exhiben marcadores emocionales
        (miedo, amenaza, imperativo) cuantitativamente distintos del resto del corpus?
  \item ¿Puede el análisis computacional aportar evidencia sobre la organización interna
        del códice --- por ejemplo, si la colocación del retrato del diablo junto a la
        Ciudad Celestial responde a un contraste deliberado que se refleja también en la
        selección textual circundante?
\end{enumerate}

\section{Metodología propuesta}

\subsection{Corpus y digitalización}

No parto de las imágenes del manuscrito: la transcripción directa (HTR, \emph{Handwritten
Text Recognition} sobre minúscula carolina tardía) es un proyecto en sí mismo. La estrategia
práctica es ensamblar el corpus desde ediciones digitales existentes de los textos que el
códice contiene --- la Vulgata, Josefo, Isidoro y Cosmas están disponibles en repositorios
como Perseus y la Latin Library --- y usar la digitalización de la Biblioteca Nacional de
Suecia para verificar la composición y las secciones exclusivas (conjuros, calendario).

\subsection{Modelos: NLP para una lengua sin hablantes}

El reto central es que el corpus está en latín medieval, una lengua de ``bajos recursos''
para el NLP moderno: no hay léxicos de sentimiento nativos ni corpus anotados de emociones.
El plan combina tres niveles:

\begin{enumerate}
  \item \textbf{Base contextual}: Latin BERT (Bamman y Burns, 2020), el primer modelo de
        lenguaje contextual para latín, entrenado con 642.7 millones de tokens que abarcan
        22 siglos de producción escrita. Provee representaciones vectoriales sensibles al
        contexto para cada palabra del corpus.
  \item \textbf{Proyección de léxicos}: transferir léxicos de emociones modernos (como el
        NRC Emotion Lexicon) al latín vía los espacios de embeddings, con validación manual
        de las entradas más frecuentes --- la traducción automática de léxicos es ruidosa y
        el vocabulario religioso medieval tiene cargas semánticas propias
        (\emph{timor Domini}, el temor de Dios, no es simplemente ``miedo'').
  \item \textbf{Clasificación supervisada ligera}: anotar manualmente una muestra
        estratificada de pasajes (con criterios explícitos y verificación de acuerdo entre
        anotadores) para afinar un clasificador sobre las representaciones de Latin BERT,
        siguiendo la estrategia que ha funcionado en otras lenguas históricas como el chino
        clásico.
\end{enumerate}

\subsection{Análisis}

Con el corpus etiquetado, el análisis es estadístico: distribución de polaridad y emociones
por obra y por género, series ``narrativas'' de emoción a lo largo de cada libro (el arco
emocional de los evangelios frente al del Apocalipsis, por ejemplo), y contraste de
hipótesis sobre los exorcismos frente a secciones de control. La visualización final ---
un mapa emocional navegable del códice --- se publicará como demo interactiva en mi sitio,
como hice con mi anterior proyecto de detección de paráfrasis.

\section{Retos y honestidad metodológica}

Este proyecto exige cautelas que conviene declarar desde el diseño. Primera: los conceptos
emocionales son culturales; proyectar categorías del siglo XXI sobre textos del XIII es un
acto interpretativo que debe hacerse explícito, no esconderse tras la aparente neutralidad
del número. Segunda: el latín de la Vulgata es traducción de originales hebreos y griegos
--- el sentimiento medido es el del texto latino que el escriba copió, no ``el de la
Biblia''. Tercera: la validación con especialistas en filología medieval no es opcional;
un clasificador puede ser estadísticamente sólido y filológicamente absurdo. El valor del
proyecto está precisamente en el diálogo entre la escala computacional (imposible para la
lectura humana) y el juicio experto (imposible para la máquina).

\section{Resultados esperados}

Espero producir: (1) un corpus latino estructurado y alineado con la composición del códice,
reutilizable por otros investigadores; (2) un léxico de sentimientos para latín medieval
validado parcialmente a mano; (3) el análisis emocional comparativo entre géneros del códice;
y (4) la visualización interactiva pública. Más allá de los resultados concretos, la apuesta
de fondo es metodológica: demostrar que las humanidades digitales pueden tratar objetos
patrimoniales complejos con el mismo rigor técnico que cualquier corpus moderno --- y que un
manuscrito del siglo XIII envuelto en leyendas diabólicas es, también, un dataset.

\section*{Referencias principales}

\begin{itemize}
  \item Bamman, D. y Burns, P.~J., \emph{Latin BERT: A Contextual Language Model for Classical
        Philology} (2020). \url{https://arxiv.org/abs/2009.10053}
  \item Mohammad, S., \emph{Sentiment Analysis: Automatically Detecting Valence, Emotions, and
        Other Affectual States from Text} (2020). \url{https://arxiv.org/abs/2005.11882}
  \item Biblioteca Nacional de Suecia, \emph{The Codex Gigas}.
        \url{https://www.kb.se/eng/the-codex-gigas/}
  \item Riemenschneider, F. y Frank, A., \emph{Exploring Large Language Models for Classical
        Philology} (2023). \url{https://arxiv.org/abs/2305.13698}
\end{itemize}

\end{document}
